\section{Tabellen}
\label{sec:tabellen}

Hier gibt es ein paar Tabellen. Um sie in VS Code schnell schreiben zu können, bietet der \LaTeX-Workshop einige Autocomplete-Funktionen wie \textit{BTA}:

\begin{table}[ht]
    \centering
    \begin{tabular}{|c|cc|}
        \hline
         & A & B \\
        \hline
        C & 1 & 0 \\
        D & 0 & 0 \\
        \hline
    \end{tabular}
    \caption{Eine einfache Übergangstabelle}
    \label{tab:simpleTable}
\end{table}

Bei größeren Tabellen eignet sich auch \textit{tabularx}, da dies automatische Zeilenumbrüche ermöglicht.

\begin{table}[ht]
    \centering
    \begin{tabularx}{0.9\textwidth}{|X|X|}
        \hline
        \LaTeX & \LaTeX ist ein plattformunabhängiges und freies Softwarepaket, das die Benutzung des Textsatzsystems TeX mit Hilfe von Makros vereinfacht.\\
        \hline
        Lorem Ipsum & Lorem ipsum dolor sit amet, consetetur sadipscing elitr, sed diam nonumy eirmod tempor invidunt ut labore et dolore magna aliquyam erat, sed diam voluptua. \\
        \hline
    \end{tabularx}
    \caption{Eine Tabelle die tabularx verwendet}
    \label{tab:tabularx}
\end{table}

Mit \textit{multirow} können auch explizit Zellen einer Tabelle verbunden werden:

\begin{table}[ht]
    \centering
    \begin{tabular}{|c|c|c|}
        \hline
        \multirow{2}{*}{1} & 2 & 3 \\
         & 4 & 5 \\\hline
        6 & 7 & 8 \\\hline
    \end{tabular}
    \caption{Tabelle mit multirow}
    \label{tab:multirow}
\end{table}

Weitere Informationen findest du unter \href{https://ctan.org/pkg/multirow}{CTAN multirow} und \href{https://ctan.org/pkg/tabularx}{CTAN tabularx}.